\documentclass[12pt, a4paper]{article}
\usepackage[utf8]{inputenc}
\usepackage[T1]{fontenc}
\usepackage[spanish]{babel}
\usepackage{graphicx}
\usepackage{booktabs}
\usepackage{geometry}
\usepackage{fancyhdr}
\usepackage{hyperref}

\geometry{margin=2.5cm}
\pagestyle{fancy}
\fancyhf{}
\chead{GLUD Fried Chicken - Taller de Metricas de Software}
\lhead{}
\rhead{}
\cfoot{\thepage}

\begin{document}

\begin{titlepage}
  \centering
  \vspace*{2cm}

  {\Huge\bfseries GLUD Fried Chicken (GFC)\par}
  \vspace{1cm}
  {\Large Taller de Metricas de Software\par}
  \vspace{0.5cm}
  {\large Diseno de Software\par}

  \vspace{2cm}
  \begin{tabular}{rl}
    \textbf{Nombre} & \textbf{Codigo} \\
    \hline
    Daniel Mateo Ballesteros Molina & 20221020102 \\
    Jan Henrik Sanchez Jerez & 20231020130 \\
    Juan Carlos Quintero Rubiano & 20232020172 \\
    Juan Diego Peralta Barbosa & 20231020146 \\
  \end{tabular}

  \vspace{3cm}
  {\large \today\par}
\end{titlepage}

\tableofcontents
\newpage

\section{Introduccion}
Completar con la introduccion del taller y el objetivo de la aplicacion.

\section{Descripcion de la aplicacion}
Describir la aplicacion CRUD para la gestion de comidas rapidas: entidades,
funcionalidades y decisiones de diseno tomadas por el grupo.

\section{Metricas con Source Monitor}
\subsection{Instalacion}
Pasos para instalar y configurar Source Monitor sobre el codigo fuente Java.

\subsection{Metricas manejadas por la herramienta}
Explicar el tipo de metricas que calcula Source Monitor (complejidad
ciclomatica, lineas de codigo, sentencias por metodo, profundidad de bloques,
etc.) y su significado.

\subsection{Resultados obtenidos}
Incluir la tabla de resultados de las metricas sobre la aplicacion:
\begin{table}[h]
  \centering
  \begin{tabular}{lcc}
    \toprule
    \textbf{Metrica} & \textbf{Valor} & \textbf{Rango aceptable} \\
    \midrule
    Lineas de Codigo (M0) & & \\
    Sentencias (M1) & & \\
    Metodos por Clase (M6) & & \\
    Sentencias por Metodo (M7) & & \\
    Complejidad Maxima (M10) & & \\
    Profundidad Promedio (M13) & & \\
    \bottomrule
  \end{tabular}
  \caption{Metricas obtenidas con Source Monitor}
\end{table}

\begin{figure}[h]
  \centering
  \includegraphics[width=0.7\textwidth]{kiviat.png}
  \caption{Diagrama Kiviat generado automaticamente por el pipeline}
  \label{fig:kiviat}
\end{figure}

\subsection{Interpretacion de resultados}
Analizar los valores obtenidos: que areas del codigo estan bien, cuales
podrian mejorarse, y como se comparan con los rangos aceptables.

\section{Metricas con SonarQube}
SonarQube complementa a Source Monitor incorporando metricas de calidad
orientadas a la norma ISO/IEC 25010 (mantenibilidad, fiabilidad, seguridad
y eficiencia), ademas de indicadores de duplicacion de codigo, code smells,
cobertura de pruebas y deuda tecnica. Mientras que Source Monitor mide
metrica interna del codigo fuente, SonarQube las interpreta dentro de
atributos de calidad externos.

\subsection{Metricas manejadas por SonarQube}
\begin{itemize}
  \item Complejidad ciclomatica y distribucion por metodo.
  \item Cobertura de codigo por lineas y por ramas.
  \item Code Smells: malos olores de codigo que afectan la mantenibilidad.
  \item Duplicacion: porcentaje de lineas duplicadas.
  \item Deuda tecnica y Maintainability Rating (A - E) segun ISO 25010.
  \item Vulnerabilidades y reglas de seguridad (fiabilidad/seguridad).
\end{itemize}

\subsection{Resultados obtenidos con SonarQube}
\begin{table}[h]
  \centering
  \begin{tabular}{lcc}
    \toprule
    \textbf{Metrica} & \textbf{Valor} & \textbf{Rango aceptable} \\
    \midrule
    Maintainability Rating (A-E) & & \\
    Reliability Rating (A-E) & & \\
    Security Rating (A-E) & & \\
    Complejidad Ciclomatica & & \\
    Cobertura (\%) & & \\
    Lineas duplicadas (\%) & & \\
    Deuda tecnica (horas) & & \\
    Code Smells & & \\
    \bottomrule
  \end{tabular}
  \caption{Metricas obtenidas con SonarQube}
\end{table}

\subsection{Interpretacion de resultados (SonarQube)}
Relacionar los ratings de SonarQube con las caracteristicas de calidad de la
ISO 25010 y comparar con las metricas de Source Monitor para el mismo codigo.

\section{Comparacion Source Monitor vs SonarQube}
Comparar el tipo de metricas que maneja cada herramienta: alcance,
granularidad, relacion con las normas de calidad y que aspectos del codigo
cubre cada una.

\section{Conclusiones}
Cerrar con las conclusiones del grupo sobre la experiencia de aplicar
metricas de software (Source Monitor y SonarQube) y la calidad del codigo
desarrollado.

\end{document}
